\documentclass[12pt,a4paper]{article}

\usepackage[utf8]{inputenc}
\usepackage[T2A]{fontenc}
\usepackage[russian]{babel}
\usepackage{amsmath, amssymb, amsthm}
\usepackage{geometry}
\usepackage{tikz}
\usepackage{float}

\geometry{top=2cm, bottom=2cm, left=2.5cm, right=1.5cm}

\newtheorem{theorem}{Теорема}
\newtheorem{definition}{Определение}
\newtheorem{lemma}{Лемма}

\title{Билет №88 \\ \small Последовательное и параллельное соединение систем. (СпБПУ)}
\author{Сериков Владислав Александрович}
\date{16 мая 2026}

\begin{document}

\maketitle
\tableofcontents
\newpage

\section{Введение}

Для анализа и синтеза сложных систем автоматического управления (САУ) и обработки сигналов применяется декомпозиция: сложная система разбивается на простые элементарные звенья. Связи между этими звеньями определяют топологию системы. Основными видами соединений являются последовательное, параллельное и соединение с обратной связью (в данном конспекте подробно рассматриваются первые два).

\begin{definition}
    \textbf{Линейная стационарная система (ЛСС)} — это система, для которой выполняется принцип суперпозиции (линейность) и чьи характеристики не меняются с течением времени (стационарность).
\end{definition}

В теории управления ЛСС могут быть описаны двумя основными способами:
\begin{enumerate}
    \item \textbf{Передаточная функция} $W(s)$ (в частотной области через преобразование Лапласа).
    \item \textbf{Уравнения в пространстве состояний} (временная область) вида:
    \begin{equation}
        \begin{cases}
            \dot{x}(t) = Ax(t) + Bu(t) \\
            y(t) = Cx(t) + Du(t)
        \end{cases}
    \end{equation}
    где $x(t)$ — вектор состояний, $u(t)$ — вектор входа, $y(t)$ — вектор выхода, а $A, B, C, D$ — матрицы системы.
\end{enumerate}

\section{Последовательное соединение систем}

\begin{definition}
    \textbf{Последовательным (каскадным) соединением} называется такое соединение, при котором выход предыдущего звена подается на вход последующего без разветвлений и внешних воздействий.
\end{definition}

\begin{figure}[H]
    \centering
    \begin{tikzpicture}[auto, node distance=2.5cm, >=latex]
        \node[coordinate] (input) {};
        \node[draw, rectangle, minimum height=1cm, minimum width=1.5cm, right of=input, node distance=2cm] (sys1) {$W_1(s)$};
        \node[draw, rectangle, minimum height=1cm, minimum width=1.5cm, right of=sys1, node distance=3cm] (sys2) {$W_2(s)$};
        \node[coordinate, right of=sys2, node distance=2.5cm] (output) {};

        \draw[->] (input) -- node {$u(t)$} (sys1);
        \draw[->] (sys1) -- node {$y_1(t) = u_2(t)$} (sys2);
        \draw[->] (sys2) -- node {$y(t)$} (output);
    \end{tikzpicture}
    \caption{Структурная схема последовательного соединения}
\end{figure}

\subsection{Описание через передаточные функции}

\begin{theorem}[Эквивалентная передаточная функция последовательного соединения]
    Передаточная функция $W_{eq}(s)$ системы, состоящей из $n$ последовательно соединенных звеньев с передаточными функциями $W_i(s)$, равна произведению передаточных функций этих звеньев:
    \begin{equation}
        W_{eq}(s) = \prod_{i=1}^{n} W_i(s)
    \end{equation}
\end{theorem}

\begin{proof}
    Рассмотрим систему из двух последовательно соединенных звеньев. Пусть $U(s)$ — изображение по Лапласу входного сигнала, $Y_1(s)$ — изображение промежуточного сигнала (выхода первого звена), $Y(s)$ — изображение выходного сигнала всей системы.

    По определению передаточной функции для каждого звена:
    \begin{equation}
        Y_1(s) = W_1(s)U(s)
    \end{equation}
    \begin{equation}
        Y(s) = W_2(s)Y_1(s)
    \end{equation}

    Подставим выражение для $Y_1(s)$ во второе уравнение:
    \begin{equation}
        Y(s) = W_2(s) \big(W_1(s)U(s)\big) = \big(W_2(s)W_1(s)\big) U(s)
    \end{equation}

    Отсюда эквивалентная передаточная функция:
    \begin{equation}
        W_{eq}(s) = \frac{Y(s)}{U(s)} = W_2(s)W_1(s) = W_1(s)W_2(s)
    \end{equation}
    Для $n$ звеньев доказательство проводится методом математической индукции. Теорема доказана.
\end{proof}

\subsection{Описание в пространстве состояний}

\begin{theorem}[Последовательное соединение в пространстве состояний]
    Пусть заданы две системы в пространстве состояний:
    Система 1: $\dot{x}_1 = A_1 x_1 + B_1 u$, \quad $y_1 = C_1 x_1 + D_1 u$ \\
    Система 2: $\dot{x}_2 = A_2 x_2 + B_2 u_2$, \quad $y = C_2 x_2 + D_2 u_2$

    При последовательном соединении ($u_2 = y_1$) эквивалентная система с вектором состояния $x = \begin{bmatrix} x_1 \\ x_2 \end{bmatrix}$ имеет матрицы:
    \begin{equation}
        A_{eq} = \begin{bmatrix} A_1 & 0 \\ B_2 C_1 & A_2 \end{bmatrix}, \quad
        B_{eq} = \begin{bmatrix} B_1 \\ B_2 D_1 \end{bmatrix}, \quad
        C_{eq} = \begin{bmatrix} D_2 C_1 & C_2 \end{bmatrix}, \quad
        D_{eq} = D_2 D_1
    \end{equation}
\end{theorem}

\begin{proof}
    Подставим выход первой системы $y_1$ на вход второй $u_2$:
    \begin{equation}
        \dot{x}_2 = A_2 x_2 + B_2(C_1 x_1 + D_1 u) = B_2 C_1 x_1 + A_2 x_2 + B_2 D_1 u
    \end{equation}
    \begin{equation}
        y = C_2 x_2 + D_2(C_1 x_1 + D_1 u) = D_2 C_1 x_1 + C_2 x_2 + D_2 D_1 u
    \end{equation}

    Запишем дифференциальные уравнения для объединенного вектора состояния $x = \begin{bmatrix} x_1 \\ x_2 \end{bmatrix}$:
    \begin{equation}
        \begin{bmatrix} \dot{x}_1 \\ \dot{x}_2 \end{bmatrix} =
        \begin{bmatrix} A_1 x_1 + B_1 u \\ B_2 C_1 x_1 + A_2 x_2 + B_2 D_1 u \end{bmatrix} =
        \begin{bmatrix} A_1 & 0 \\ B_2 C_1 & A_2 \end{bmatrix} \begin{bmatrix} x_1 \\ x_2 \end{bmatrix} +
        \begin{bmatrix} B_1 \\ B_2 D_1 \end{bmatrix} u
    \end{equation}
    Для выходного уравнения:
    \begin{equation}
        y = \begin{bmatrix} D_2 C_1 & C_2 \end{bmatrix} \begin{bmatrix} x_1 \\ x_2 \end{bmatrix} + D_2 D_1 u
    \end{equation}
    Сравнивая с классической формой пространства состояний, получаем требуемые матрицы. Теорема доказана.
\end{proof}

\section{Параллельное соединение систем}

\begin{definition}
    \textbf{Параллельным соединением (согласно-параллельным)} называется такое соединение, при котором один и тот же входной сигнал подается на входы всех звеньев, а их выходные сигналы алгебраически суммируются.
\end{definition}

\begin{figure}[H]
    \centering
    \begin{tikzpicture}[auto, node distance=2cm, >=latex]
        \node[coordinate] (input) {};
        \node[fill, circle, inner sep=1.5pt, right of=input, node distance=1.5cm] (branch) {};
        \node[draw, rectangle, minimum height=1cm, minimum width=1.5cm, right of=branch, node distance=2.5cm, yshift=1cm] (sys1) {$W_1(s)$};
        \node[draw, rectangle, minimum height=1cm, minimum width=1.5cm, right of=branch, node distance=2.5cm, yshift=-1cm] (sys2) {$W_2(s)$};
        \node[circle, draw, right of=branch, node distance=5.5cm, inner sep=0pt, minimum size=0.5cm] (sum) {$+$};
        \node[coordinate, right of=sum, node distance=1.5cm] (output) {};

        \draw[->] (input) -- node {$u(t)$} (branch);
        \draw[->] (branch) |- (sys1);
        \draw[->] (branch) |- (sys2);
        \draw[->] (sys1) -| node[pos=0.9, right] {$+$} (sum);
        \draw[->] (sys2) -| node[pos=0.9, right] {$+$} (sum);
        \draw[->] (sum) -- node {$y(t)$} (output);
    \end{tikzpicture}
    \caption{Структурная схема параллельного соединения}
\end{figure}

\subsection{Описание через передаточные функции}

\begin{theorem}[Эквивалентная передаточная функция параллельного соединения]
    Передаточная функция $W_{eq}(s)$ системы, состоящей из $n$ параллельно соединенных звеньев с передаточными функциями $W_i(s)$, равна алгебраической сумме передаточных функций этих звеньев:
    \begin{equation}
        W_{eq}(s) = \sum_{i=1}^{n} W_i(s)
    \end{equation}
\end{theorem}

\begin{proof}
    Рассмотрим два звена. Изображение входного сигнала $U(s)$ является общим.
    Выходы первого и второго звеньев:
    \begin{equation}
        Y_1(s) = W_1(s)U(s), \quad Y_2(s) = W_2(s)U(s)
    \end{equation}
    Выход системы равен сумме выходов отдельных звеньев (свойство линейности преобразования Лапласа):
    \begin{equation}
        Y(s) = Y_1(s) + Y_2(s) = W_1(s)U(s) + W_2(s)U(s) = \big(W_1(s) + W_2(s)\big)U(s)
    \end{equation}
    Отсюда:
    \begin{equation}
        W_{eq}(s) = \frac{Y(s)}{U(s)} = W_1(s) + W_2(s)
    \end{equation}
    Для $n$ звеньев результат обобщается по индукции. Теорема доказана.
\end{proof}

\subsection{Описание в пространстве состояний}

\begin{theorem}[Параллельное соединение в пространстве состояний]
    Для двух систем, заданных в пространстве состояний (матрицы $A_1, B_1, C_1, D_1$ и $A_2, B_2, C_2, D_2$), параллельное соединение ($u_1 = u_2 = u$, $y = y_1 + y_2$) задает эквивалентную систему с вектором состояния $x = \begin{bmatrix} x_1 \\ x_2 \end{bmatrix}$ и матрицами:
    \begin{equation}
        A_{eq} = \begin{bmatrix} A_1 & 0 \\ 0 & A_2 \end{bmatrix}, \quad
        B_{eq} = \begin{bmatrix} B_1 \\ B_2 \end{bmatrix}, \quad
        C_{eq} = \begin{bmatrix} C_1 & C_2 \end{bmatrix}, \quad
        D_{eq} = D_1 + D_2
    \end{equation}
\end{theorem}

\begin{proof}
    Запишем исходные уравнения для параллельного соединения с общим входом $u$:
    \begin{equation}
        \dot{x}_1 = A_1 x_1 + B_1 u, \quad \dot{x}_2 = A_2 x_2 + B_2 u
    \end{equation}
    Объединяя в одну матричную систему:
    \begin{equation}
        \begin{bmatrix} \dot{x}_1 \\ \dot{x}_2 \end{bmatrix} =
        \begin{bmatrix} A_1 & 0 \\ 0 & A_2 \end{bmatrix} \begin{bmatrix} x_1 \\ x_2 \end{bmatrix} +
        \begin{bmatrix} B_1 \\ B_2 \end{bmatrix} u
    \end{equation}
    Выход системы равен сумме выходов:
    \begin{equation}
        y = y_1 + y_2 = (C_1 x_1 + D_1 u) + (C_2 x_2 + D_2 u) =
        \begin{bmatrix} C_1 & C_2 \end{bmatrix} \begin{bmatrix} x_1 \\ x_2 \end{bmatrix} + (D_1 + D_2)u
    \end{equation}
    Матрицы соответствуют заявленным в теореме. Доказательство завершено.
\end{proof}

\section{Алгоритм эквивалентирования (свертки) структурных схем}

При расчете сложных линейных систем часто требуется получить передаточную функцию от одного входа к одному выходу. Для этого применяется алгоритм свертки графов (структурных схем).

\vspace{0.5cm}
\textbf{Алгоритм свертки структурной схемы:}
\begin{enumerate}
    \item \textbf{Шаг 1.} Выделить в схеме участки с чистым последовательным соединением (отсутствуют точки ветвления между звеньями). Заменить каждый участок на эквивалентное звено, умножив их передаточные функции.
    \item \textbf{Шаг 2.} Выделить участки с чистым параллельным соединением (звенья имеют общую точку начала и общую точку суммирования, без дополнительных входов/выходов). Заменить их на эквивалентное звено, сложив передаточные функции.
    \item \textbf{Шаг 3.} (Для полноты). Выделить звенья, охваченные локальной обратной связью, и свернуть их по формуле $W_{eq} = \frac{W}{1 \pm W \cdot W_{oc}}$.
    \item \textbf{Шаг 4.} Повторять Шаги 1-3 до тех пор, пока система не преобразуется в одно эквивалентное звено.
\end{enumerate}

\subsection{Минимальный пример применения алгоритма}

Пусть заданы два апериодических (инерционных) звена первого порядка:
\begin{equation}
    W_1(s) = \frac{1}{s+1}, \quad W_2(s) = \frac{2}{s+2}
\end{equation}

\textbf{Случай А: Последовательное соединение (Шаг 1)}
Требуется найти эквивалентную систему.
Применяем Теорему 1:
\begin{equation}
    W_{ser}(s) = W_1(s) \cdot W_2(s) = \frac{1}{s+1} \cdot \frac{2}{s+2} = \frac{2}{s^2 + 3s + 2}
\end{equation}
Получили колебательное/апериодическое звено второго порядка.

\textbf{Случай Б: Параллельное соединение (Шаг 2)}
Применяем Теорему 3:
\begin{equation}
    W_{par}(s) = W_1(s) + W_2(s) = \frac{1}{s+1} + \frac{2}{s+2} = \frac{(s+2) + 2(s+1)}{(s+1)(s+2)} = \frac{3s+4}{s^2 + 3s + 2}
\end{equation}
В параллельном соединении, как видно из результата, в числителе эквивалентной передаточной функции возникают нули (корни числителя), влияющие на переходный процесс (в данном случае нуль $s = -4/3$). Это важное отличие параллельного соединения от последовательного, при котором нули образуются только если они присутствовали в исходных звеньях.

\end{document}
